% from ega4-16-fr.tex
\oldpage[IV-4]{9}
locale, et sans que $f$ soit injectif. On en a un exemple en prenant pour $Y$
une somme de préschémas $Y_\lambda$ tous isomorphes à
$\operatorname{Spec}(\mathcal{O}_x)$, où $x\in X$, et pour $f$ le morphisme
égal au morphisme canonique dans chacun des $Y_\lambda$.
\end{enumerate}
\end{remarks}

\subsection{Propriétés fonctorielles des invariants normaux d'une immersion}
\EGAReaderTarget{ega4.c04.subsec.16-2}{IV.16.2}

\begin{env}[16.2.1]
\label{IV.16.2.1}
Soient $f=(\psi,\theta):Y\to X$,
$f'=(\psi',\theta'):Y'\to X'$ deux morphismes d'espaces annelés tels que les
homomorphismes $\theta^\#$ et $\theta'^\#$ soient surjectifs; considérons un
diagramme commutatif de morphismes d'espaces annelés
\[
  \xymatrix{
    Y \ar[r]^f & X \\
    Y' \ar[r]_{f'} \ar[u]^u & X' \ar[u]_v
  }
  \tag{16.2.1.1}
  \label{IV.16.2.1.1}
\]
Posons $u=(\rho,\lambda)$, $v=(\sigma,\mu)$. On a
$\rho^*(\psi^*(\mathcal{O}_X))=
\psi'^*(\sigma^*(\mathcal{O}_X))$ et on a par suite un diagramme commutatif
d'homomorphismes de faisceaux d'anneaux sur $Y'$
\[
  \xymatrix{
    \rho^*(\psi^*(\mathcal{O}_X))
      =\psi'^*(\sigma^*(\mathcal{O}_X))
      \ar[r]^-{\psi'^*(\mu^\#)}
      \ar[d]_{\rho^*(\theta^\#)}
      & \psi'^*(\mathcal{O}_{X'}) \ar[d]^{\theta'^\#} \\
    \rho^*(\mathcal{O}_Y) \ar[r]_-{\lambda^\#}
      & \mathcal{O}_{Y'}
  }
\]
d'où l'on conclut, si $\mathscr{I}$ et $\mathscr{I}'$ sont les noyaux de
$\theta^\#$ et $\theta'^\#$, que l'on a
$\psi'^*(\mu^\#)(\rho^*(\mathscr{I}))\subset\mathscr{I}'$, compte tenu de
l'exactitude du foncteur $\rho^*$. On en déduit aussitôt que pour tout entier
$n$,
$\psi'^*(\mu^\#)(\rho^*(\mathscr{I}^n))\subset\mathscr{I}'^n$, ce qui montre
que $\psi'^*(\mu^\#)$ définit, par passage aux quotients, un homomorphisme de
faisceaux d'anneaux
\[
  \nu_n:
  \rho^*(\psi^*(\mathcal{O}_X)/\mathscr{I}^{n+1})
  \longrightarrow
  \psi'^*(\mathcal{O}_{X'})/\mathscr{I}'^{n+1}
  \tag{16.2.1.2}
  \label{IV.16.2.1.2}
\]
et par suite un morphisme d'espaces annelés
$w_n=(\rho,\nu_n):Y'^{(n)}\to Y^{(n)}$ (qui, pour $n=0$, n'est autre que
$u$). Il résulte aussitôt de cette définition que les diagrammes
\[
  \xymatrix@R=1pc{
    Y^{(n)} \ar[r]^-{h_{mn}} & Y^{(m)} \ar[r]^-{h_m} & X \\
    & & & (n\leq m) \\
    Y'^{(n)} \ar[r]_-{h'_{mn}} \ar[uu]^-{w_n}
      & Y'^{(m)} \ar[r]_-{h'_m} \ar[uu]^-{w_m}
      & X' \ar[uu]_-v
  }
\]
(où les flèches horizontales sont les morphismes canoniques (16.1.2)) sont
commutatifs.

Par passage aux quotients à partir des homomorphismes (16.2.1.2), et en tenant
